% !PRINTED 433
\weberpage{433}
%<<CONT
or
\weq{14}{\beta\,\omega^2 + (\beta\,c'+\beta\,c+\alpha\,p-\delta\,p)\,\omega + \beta\,c\,c'+\alpha\,c\,p-\delta\,c'\,p-\gamma\,p^2 = 0,}
whence by comparison with (4):
\weq{15}{\begin{aligned}
\beta &= A\,x,\\
\beta\,c'+\beta\,c+\alpha\,p-\delta\,p &= B\,x,\\
\beta\,c\,c'+\alpha\,c\,p-\delta\,c'\,p-\gamma\,p^2 &= C\,x.
\end{aligned}}
and if one again sets
\weq{16}{\beta\,c'-\beta\,c+\alpha\,p+\delta\,p = y}
then:
\weq{17}{\begin{aligned}
2\,(\beta\,c'+\alpha\,p) &= B\,x+y,\\
2\,(\beta\,c-\delta\,p) &= B\,x-y.
\end{aligned}}
From this
\weq{18}{y^2+m\,x^2 = 4\,p^2.}
If $x$ is divisible by $p$, then $y$ must also be divisible by $p$, and by cancellation of the factor $p^2$ we come back to equation (11), which, as we have seen, is solvable only in the cases $m=3$ and $m=4$. We therefore assume further that $x$ is indivisible by $p$.

If $m$ is divisible by $p$, then $y$ is also divisible by $p$, and $\beta$ is, by (15), not divisible by $p$. But then it follows from (16): $c\equiv c' \pmod{p}$, and both therefore correspond to the same root of (1).

So let $p$ now divide neither $x$ nor $m$. Then it does not divide $y$ either, and by (17) $c$ and $c'$ are determined from the congruences
\weq{19}{\beta\,c' \equiv \frac{B\,x+y}{2}, \quad \beta\,c \equiv \frac{B\,x-y}{2} \pmod{p}}
and once $c,\,c'$ are determined, $\alpha$ and $\delta$ result from (17) as \gesperrt{integers}. With this the first two equations (15) are satisfied, and $\gamma$ is obtained from the last equation (15) likewise as an integer.

Indeed we have
\weq{20}{\gamma\,p^2 = \beta\,c\,c'+(\alpha\,c-\delta\,c')\,p-C\,x \equiv 0 \pmod{p^2}.}
For from (17) and (15) we have
\[
\begin{aligned}
\beta^2\,c\,c'+\beta\,(\alpha\,c-\delta\,c')\,p-\alpha\,\delta\,p^2-A\,C\,x^2
&= \beta\,[\beta\,c\,c'+(\alpha\,c-\delta\,c')\,p-C\,x]-\alpha\,\delta\,p^2\\
&= -\frac{y^2+m\,x^2}{4} = -p^2;
\end{aligned}
\]
%CONT>>
